\section{Contexto}
\label{sec:contexto}

El problema abordado es la ineficiencia de gestionar manualmente transacciones bancarias debido a la alta variabilidad de sus descripciones. Como se observa en los datos del proyecto, un mismo gasto puede aparecer registrado de formas muy distintas: desde términos breves y simples hasta oraciones complejas y formales.

Esta heterogeneidad, sumada al volumen de movimientos, impide el uso de reglas fijas o filtros de texto simples, que fallan ante cualquier cambio en la redacción. Por ello, es necesario un sistema basado en Procesamiento de Lenguaje Natural (PLN) capaz de interpretar el significado de estas descripciones para clasificarlas automáticamente en categorías específicas (\textit{Area}). 

Para validar la eficacia de esta clasificación, se utilizan las siguientes métricas:

\begin{itemize}
    \item \textbf{Accuracy (Exactitud):} Mide el porcentaje total de transacciones que el modelo ha clasificado en la categoría correcta.
    \item \textbf{F1-Score:} Evalúa el equilibrio entre la precisión y la capacidad de recuperación del modelo. Es fundamental para asegurar que el sistema clasifique bien tanto las categorías muy comunes (como \textit{Food}) como aquellas con menos ejemplos en el dataset.
\end{itemize}

El objetivo es transformar este flujo de descripciones heterogéneas en información organizada y etiquetada de forma automática.